\section{The exact noisy endpoint}\label{sec:noisy-exact}
\begin{theorem}\label{thm:noisy-exact}
Fix \(B>0\) and \(E=\lfloor e^{Bn^{1/3}}\rfloor\). At
\(K=\kappa_{\rm noise}n^{-5/6}e^{(C+2B)n^{1/3}}\), the probability
of recovering all binary faces under budget \(E\) tends to zero.
\end{theorem}

\begin{proof}
We prove failure at the exact noisy endpoint by applying the second moment
method to a weighted count of confusable opposite-face pairs of codimension
near the optimal value. The main idea is to split the distance-histogram rows
into dense rows, where we enforce an exact \(\ell^1\) constraint, and nondense
rows, where we exponentially tilt the counts and truncate to a narrow window
around the tilted mean so that every positively weighted pair is guaranteed to
be confusable.

Let \(r\) be the even codimension nearest to \(\rho_*n\), and consider its
\(N=\binom nr2^{r-1}\) opposite-face candidates.
Put \(s=n^{1/3}\), \(W=n^{2/3}\), \(v=r/(4n)\),
\(R=\sqrt{2vC}\), \(\tau=RW/(2E)\), and \(\gamma=1/10\).
As before, for each \(j<r/2\), write \(S_j=H_j-H_{r-j}\) for the signed
difference between the landmark counts at distance \(j\) and distance \(r-j\),
and let \(\lambda_j=K\prob(\Bin(r,1/2)=j)\) be their common expected count.
We call a row dense if \(\lambda_j\ge E^2e^{4n^\gamma}\).
There are \(m=RW+O(sn^\gamma)\) dense rows. Let \(T\) denote the
sum of \(|S_j|\) over all nondense rows.

\step{A weight supported at the exact budget.}
To control the nondense rows, we tilt each nondense row by the factor
\(e^{-\tau |S_j|}\) while preserving independence across rows. Under this
tilted law, we have \(\ex T=O(q_0/\tau)\) and
\(\operatorname{Var}T=O(q_0/\tau^2)\), where \(q_0=O(sn^\gamma)\) bounds the
number of nondense rows satisfying \(\lambda_j\tau^2\ge n^{-3}\). The bound on
the size of this shoulder of nondense rows follows from Stirling's
approximation.

To justify these moment estimates, we partition the nondense rows into three
groups. For shoulder rows with \(\lambda_j\le\tau^{-2}\), the weight is a
decreasing function of \(|S_j|\), so monotone weighting bounds the tilted
moments by the unweighted ones, whose second moment is at most \(2\lambda_j\).
For the remaining shoulder rows, the local lower bound on point probabilities
over \(|S_j|\le\tau^{-1}\) shows that the normalizing constant is
\(\gtrsim1/(\tau\sqrt{\lambda_j})\); combining this with the maximum point
probability bounds the weighted first and second absolute moments by
\(O(\tau^{-1})\) and \(O(\tau^{-2})\) after normalization. Finally, the
remaining light rows have total tilted first and second absolute moments
\(o(\tau^{-1})\) and \(o(\tau^{-2})\) by Cauchy--Schwarz and their
\(n^{-3}\) intensity cutoff.

Now let \(T_+,T_-\) be the tilted mean of \(T\) plus and minus \(\log n\)
times the uniform standard-deviation bound \(O(\sqrt{q_0}/\tau)\), clipping
\(T_-\) below at zero. Assign each opposite-face candidate the weight
\(Y=e^{\tau(T_+-T)}\) if \(T\in[T_-,T_+]\) and its dense \(\ell^1\) norm
is at most \(2E-T_+\), and assign it weight zero otherwise. Whenever a
candidate has positive weight, its total \(\ell^1\) norm across all rows is at
most \(2E\), so positive weight implies confusability. By Chebyshev's
inequality, restricting to this window preserves a \(1-o(1)\) fraction of the
tilted tail mean. Moreover, because \(T\) lies within the window below \(T_+\)
on the support of \(Y\), we have the uniform bound
\(\sup Y\le e^{O(\sqrt{q_0}\log n)}=e^{o(s)}\).

\step{The mean at the order-\(s\) scale.}
We next compute \(\log\ex Y\) up to \(o(s)\) and verify that the expected
total weight \(N\ex Y\) grows exponentially on the \(s\) scale. Uniformly over
all target vectors in the dense \(\ell^1\) ball, the point probabilities on the
dense rows equal \((1+o(1))\prod_{\rm dense}(4\pi\lambda_j)^{-1/2}\); the
Fourier estimate proved below also establishes this uniformity. By standard
composition bounds and Stirling's formula, the integer \(\ell^1\) ball of
radius \(Q\) on \(m\) rows has log volume
\(m\log(2eQ/m)+O(\log m)+o(1)\) whenever \(Q\gg m^2\). Taking
\(Q=\lfloor 2E-T_+\rfloor\), the log volume loses
\(mT_+/(2E)+O(q_0^2/W)\) relative to a ball of radius \(2E\). Because
\(|m-RW|=O(q_0)\), the factor \(e^{\tau T_+}\) in the weight cancels this
volume loss up to \(O(q_0^2/W)=o(s)\).

Write \(M_j=\ex e^{-\tau|S_j|}\). On every dense row, we have
\(M_j=(1+o(n^{-1}))/(\tau\sqrt{\pi\lambda_j})\), which follows from the
normal comparison proved below together with the fact that
\(\tau\sqrt{\lambda_j}\) is exponentially large in \(n^\gamma\).
Substituting this identity into the product of the dense atom probability, the
ball volume, and the nondense tilted expectation, and using
\(m+m\log(RW/m)=RW+O(q_0^2/W)\), gives
\[
 \log\ex Y=2E\tau+\sum_{j<r/2}\log M_j+o(s).
\]

To evaluate the sum on the right, let \(G\) be a standard normal random
variable and define \(\phi(y)=\ex e^{-y|G|}\). We show that the sum of
\(\log M_j\) differs by \(o(1)\) from the sum of
\(\log\phi(\tau\sqrt{2\lambda_j})\). For rows with
\(\lambda_j<e^{n^{1/4}}\), both absolute logarithms are
\(O(\tau\sqrt{\lambda_j})\), which is negligible even after summing over all
such rows. On the remaining rows, we use the Poisson kernel with parameter
\(e^{-\tau}\) as a probability measure on frequencies, since its integer
Fourier coefficients are \(e^{-\tau|k|}\). Complement pairing allows us to
replace the Bernoulli counts by Poisson counts uniformly across rows, with a
relative error that is exponentially small in \(n\). The marginal
characteristic function localizes to the interval
\(|\theta|\le e^{n^\gamma}/\sqrt{\lambda_j}\). On this interval, the
characteristic function has a Gaussian quadratic exponent with uniform error
\(O(e^{4n^\gamma-n^{1/4}})\), and the kernel density equals
\(\tau/[\pi(\tau^2+\theta^2)]\) up to a relative factor of
\(1+O(\tau+\theta^2)\). Extending the quadratic integral from this interval to
the whole real line costs a superexponentially small amount, which is
negligible relative to its \(e^{-O(s)}\) main term. The extended Cauchy kernel
integral evaluates to exactly \(\phi(\tau\sqrt{2\lambda_j})\). Because all
errors are \(o(n^{-1})\) per row, their total across all rows is \(o(1)\).
The dense asymptotic formula for \(M_j\) also follows from substituting
\(t=y|G|\) in the normal integral for \(\phi\).

Now reindex the rows by \(u=r/2-j\ge1\) and set \(t_u=Cs-u^2/(2vn)\).
As in Section~\ref{sec:noiseless-exact}, uniformly in the moderate window of
rows, the fair-binomial relative entropy and Stirling prefactor give
\[
 \tau\sqrt{2\lambda_j}
 \le e\sqrt{2/\pi}\,e^{(t_u+\epsilon_n)/2},
 \qquad \epsilon_n=O(n^{-2/3}),
\]
with the constant determined by \(\kappa_{\rm noise}\); outside the moderate
window, Hoeffding's inequality gives negligible total intensity.
Write \(z=t_u+\epsilon_n\). Applying Jensen's inequality gives
\[
 \log\phi(e\sqrt{2/\pi}\,e^{z/2})
 \ge -(1+z/2)\mathbf1_{z>0}-b(z),
 \quad
 b(z)=
 \begin{cases}
 (2e/\pi)e^{z/2},&z\le0,\\
 (\pi/(2e^2))e^{-z},&z>0.
 \end{cases}
\]
The first branch follows directly from Jensen's inequality and
\(\ex |G|=\sqrt{2/\pi}\). For the second branch, we rewrite the normal
integral as \(\phi(y)=\sqrt{2/\pi}\,y^{-1}\ex e^{-X^2/(2y^2)}\) with \(X\) a
mean-one exponential random variable (whose second moment is two) and then
apply Jensen's inequality. Integrating \(b\) yields
\(\int b=4e/\pi+\pi/(2e^2)<3.68\). There are \(RW+O(1)\) levels with positive
\(z\), and the sum of the half-quadratic terms over those levels is
\(a_{r/n}(C)n/2-Cs/4+O(1)\). The shoulder Riemann-sum argument from
Section~\ref{sec:noiseless-exact} gives
\(\sum b(z)=(v/R)s\int b+o(s)\). Combining these estimates with \(2E\tau=RW\)
and Stirling's approximation for \(N\), we obtain
\[
 \log(N\ex Y)\ge (C/4-3.68v/R)s+o(s)\ge cs
\]
for a positive fixed constant \(c\). Indeed, evaluating at \(\rho=1/2\) shows
that \(C>3.3\), and since \(v/R\le1/\sqrt{8C}\), the coefficient of \(s\) is
strictly positive. Rounding the maximizing codimension to the nearest even
integer changes the extensive cancellation by only \(O(1)\).

\step{The relative joint bound.}
To finish the second moment proof, we bound the joint second moment across all
pairs of opposite-face candidates and show that the relative variance tends to
zero. We start with the typical pairs of
Section~\ref{sec:noiseless-face-proof}, defined using the signed-overlap cutoff
\(n^{-1/8}\). For prescribed dense targets on both candidates, we represent the
joint weight in frequency space using Haar Fourier measure on dense rows and
the Poisson-kernel measure on nondense rows. Complement pairing, uniform
Poissonization, and projecting out distant rows of total intensity \(o(1)\)
work just as in the previous proofs. On every retained row with
\(\lambda_j\ge e^{n^{1/4}}\), we apply the neutral singleton bound
\(\lambda_j e^{-O(n^{1/12}+\log n)}\) from the strict noisy case. This
localizes the integral to the core region
\(|\theta_j|\le e^{n^\gamma}/\sqrt{\lambda_j}\) on all such rows, because the
off-core contribution is negligible even compared to the \(e^{-O(n)}\) product
of single-core integrals.

On the remaining light rows, we keep the full fundamental frequency interval.
Cross terms involving any light row are \(o(1)\) uniformly by the shell-pair
intensity bound from Section~\ref{sec:noiseless-face-proof}: within a block,
the sum of the \(\lambda_j|\theta_j|\) terms over light rows is at most
\(\pi n e^{n^{1/4}}\), while the corresponding sum on the localized rows is at
most \(n\sqrt K e^{n^\gamma}\). Dividing their product by \(K\) and allowing
for the \(e^{o(s)}\) intensity correction shows that these cross terms vanish.
The remaining light-row marginal weights have product \(1+o(1)\). On the core,
the dense target phases are also \(1+o(1)\) uniformly, because the dense
targets satisfy \(\ell^1\le2E\) and the dense frequencies obey the core cutoff.

After absorbing a uniform \(o(1)\) quartic error on the core, we scale every
localized frequency by \(\sqrt{2\lambda_j}\). Under this rescaling, each
marginal probability density on the core is proportional to a standard Gaussian
times an even decreasing function of absolute value, coming from the core
cutoff and the Poisson kernel. The normalized coordinates are therefore
independent within each marginal and centered 1-subgaussian: by monotone
weighting, their \(\cosh(tx)\) expectations are bounded above by those of the
standard Gaussian. By the signed Walsh cross-Gram estimate, the normalized
cross block \(T_0\) has Frobenius norm \(o(1)\). For independent marginal
vectors \(U,V\) drawn from the rescaled core distributions, Jensen's
inequality and Gaussian linearization give
\[
 1\le\ex e^{-U^{\mathsf T}T_0V}
 \le\det(I-T_0T_0^{\mathsf T})^{-1/2}=1+o(1).
\]
To obtain the upper estimate, condition on \(U\) and apply subgaussianity to
\(V\); then represent the resulting factor \(e^{\|T_0^{\mathsf T}U\|^2/2}\) by
an auxiliary normal linear moment and apply subgaussianity to \(U\). Thus the
mixed joint atom is \(1+o(1)\) times the product of its single-candidate
integrals, uniformly over both dense targets. Summing over both dense balls,
using the untruncated tail weights as a joint upper bound, and applying the
relative single-candidate truncation estimate handles all typical pairs.

The remaining pair classes are all negligible relative to the squared mean.
Atypical fixed-contraction pairs satisfy the same shell bounds as before, now
incurring at most an \(\exp(o(s))\) weight loss per factor against an
\(\exp(-\Omega(n^{3/4}))\) count loss. For distinct near partners, we use the
two-shift forest of Section~\ref{sec:noiseless-exact} and condition on all
landmarks outside its cells. On the support of both weights, the concatenated
folded vector has \(\ell^1\) norm at most \(4E\). Summing along a fundamental
cut therefore confines each independent cell count to an interval of length at
most \(8E+1\). Using this interval bound in place of the single-atom factor
preserves the full two-half leading rate at \(A-2B=C\), and the number of
near-partner pairs keeps the same extensive entropy margin because
\(\ex Y=e^{-a_{r/n}(C)n/2+o(n)}\). Finally, identical pairs contribute at most
\((\sup Y)N\ex Y\), which is negligible because the log mean is positive on the
\(s\) scale and \(\sup Y=e^{o(s)}\). Hence the relative variance tends to zero,
proving failure.
\end{proof}
